\section{The {\sys} System Architecture}
\label{sec:architecture}

{\sys} has two loosely coupled components.
First, the \emph{Rule Generator} runs offline, and distills the output of an LLM rewriter into verified rules and composes them in a rule set.
Second, the \emph{Rewrite Engine} does online rewriting of incoming queries with these rules and no LLM calls.
The two components share the rule set and its \emph{Fragment Index} maps fragment keys to rule identifiers.
Figure~\ref{fig:system_arch} shows the interaction among these two components.

\textbf{Design Principles:}
Three principles shape {\sys}, each serving a requirement from Section~\ref{sec:problem}.
First, {\sys} distills transformation rules, not queries. A rule refers to operator patterns, so it fires on new queries with other characteristics (R3).
Second, {\sys} verifies each rule once, before admission, not per query, so it adds no cost at query time (R1, R2).
Third, {\sys} reuses the host engine: it rewrites only the logical plan and returns SQL (R4).
The host keeps its physical optimization, parallel execution, and memory management, so the speedup of a rewrite adds additional improvements.
The LLM rewriter is pluggable: any tool that rewrites SQL (and reports the applied transformations and priorities) is applicable here.\shorten

\subsection{Rule Generator}
\label{sec:arch_generator}
The Rule Generator expands the rule set in five steps, shown on the right of Figure~\ref{fig:system_arch}.
First, it selects a sample of the workload queries $\mathcal{Q}_s \subseteq \mathcal{W}$. The sample may include queries flagged by runtime monitoring or by administrators.
Second, the LLM rewriter produces a candidate rewrite $Q''$ for each sampled query $Q$.
Third, the generator distills each pair $(Q, Q'')$ into a candidate rule.
It compares the two logical plans, uses the rewriter's explanation of its transformations, and generalizes changes beyond the source query. Repeated rewrites are simply discarded.
Fourth, an LLM implements each new rule as a C++ class with \texttt{Matches()} and \texttt{Apply()} methods.
Fifth, the rule passes the three verification stages of Definition~\ref{def:verification}: a smoke test, a formal check, and a result-equivalence test if the formal check does not succeed.
A failed stage returns a diagnostic to the LLM, which regenerates the implementation.
Fragment keys of admitted rules are added to the Fragment Index without a need for restarting the Rewrite Engine.
Section~\ref{sec:rule_generator} describes the details.

\subsection{Rewrite Engine}
\label{sec:arch_engine}
The Rewrite Engine processes each query in four steps, shown on the left of Figure~\ref{fig:system_arch}.
First, it parses the SQL query with the DuckDB parser~\cite{DuckDB} and lowers the parse tree into a logical plan of nine operator types: \textsf{Scan}, \textsf{Filter}, \textsf{Project}, \textsf{Join}, \textsf{Aggregate}, \textsf{Sort}, \textsf{Limit}, \textsf{SetOp}, and \textsf{CTE}.
All rules operate on this representation, independent of the SQL dialect.
Second, the engine extracts the fragment keys of the plan and looks them up in the Fragment Index. The result is the candidate rule set $\mathcal{R}(Q)$, and only these rules are evaluated.
Third, the engine applies the candidate rules until no rule fires or a budget is reached. Rules fire in a priority order. Section~\ref{sec:query_rewriting} describes this ordering and an optional search over alternative orders.
Fourth, we serialize the rewritten plan into the target SQL dialect.
A query that matches no rule, or cannot be parsed, passes through unaltered.
{\sys} never blocks a query and never needs an LLM at query time.\shorten

\textbf{API:} Applications use {\sys} through a small API, where the host engine gets \texttt{new\_sql} as ordinary SQL without modifications:
\begin{lstlisting}[style=ide, language=C++, morekeywords={RuleEngine, ActiveRuleIds, string}]
/* Load the verified rule set and its Fragment Index */
RuleEngine engine;
/* Extract fragment keys and look up the Fragment Index */
ActiveRuleIds rules = engine.Match(sql);
/* Apply matched rules and unparse to the target dialect */
string new_sql = engine.Rewrite(sql, rules, "PostgreSQL");
\end{lstlisting}

\begin{example}[Running Example]
Stats-CEB query~88 counts the rows of a star join between \texttt{Users}, \texttt{Comments}, and \texttt{Badges}.
The LLM rewriter decomposes this join into per-branch pre-aggregations.
The Rule Generator then distills three rules from this query/rewritten-query pair, including a rule that decomposes a \texttt{COUNT(*)} star join over any number of branches.
Query~122 joins \texttt{Users}, \texttt{Posts}, \texttt{PostHistory}, and \texttt{Badges}, which adds a branch and uses two tables not contained in query~88.
The Rewrite Engine finds the fragment key \texttt{AGG(JOIN[INNER])} in the Fragment Index and applies the same three rules to query~122 without LLM calls.
We describe how rules are distilled in Section~\ref{sec:rule_generator}, and how query~122 is rewritten in Section~\ref{sec:query_rewriting}.
\end{example}